% Candidate V1.3 extension. Do not include in the frozen headline manuscript
% until its paired robustness and nominal-noninferiority gates pass.
\section{RCS-Robust Joint Fusion--Bundle Optimization}
\label{sec:v13-rcs-robust}

Let $\bar\sigma_q$ be the scheduler-side nominal RCS and assume the declared
target-class interval
\begin{equation}
\mathcal U_q=[\underline\sigma_q,\bar\sigma_q],\qquad
\underline\sigma_q=\zeta_q\bar\sigma_q,quad 0<\zeta_q\leq1.
\end{equation}
For fixed geometry, power, interference, and waveform impairments, the
bistatic model gives
\begin{equation}
\gamma^{\rm s}_{ijq}(\sigma_q)=C_{ijq}\sigma_q,
\qquad C_{ijq}\geq0.
\end{equation}
The finite-look LLR mean gap
$L\gamma^2/(1+\gamma)$ and deflection $L\gamma^2$ are nondecreasing in
$\gamma\geq0$.  Because report success is determined by the inter-UAV link,
the fixed-bundle detector surrogate $p_{qfb}(\sigma_q)$ is nondecreasing in
$\sigma_q$.  Therefore
\begin{equation}
\min_{\sigma_q\in\mathcal U_q}p_{qfb}(\sigma_q)
=p_{qfb}(\underline\sigma_q)\equiv p^{\rm R}_{qfb}.
\label{eq:v13-endpoint}
\end{equation}

We replace the nominal reliability coefficients in the V1.2 bundle master by
$p^{\rm R}_{qfb}$ and solve
\begin{equation}
\operatorname{lexmin}\left(d^{\rm R}_{\max},\sum_qd^{\rm R}_q,
\sum_{qfb}r_{qfb}u_{qfb},\sum_{qfb}c_{qfb}u_{qfb}\right),
\end{equation}
subject to the original assignment and hard resource constraints and
\begin{align}
d^{\rm R}_q&\geq P_D^{\rm req}-\sum_{f,b}p^{\rm R}_{qfb}u_{qfb},\\
d^{\rm R}_{\max}&\geq d^{\rm R}_q.
\end{align}
Equation~\eqref{eq:v13-endpoint} makes this an exact interval-robust reduction,
not a weighted RCS heuristic.  Operationally, we scale each target's sensing
SINR by $\zeta_q$, recompute the nonlinear LLR moments, and run the unchanged
dual-guided coarse-to-fine pricing algorithm on these lower-endpoint tables.
The method has the same master dimension as V1.2 and reduces exactly to it when
$\zeta_q=1$.

The interval describes slow target-class uncertainty, not an observable RCS
realization from the current CPI.  The endpoint reduction does not cover uncertain
aspect responses that reorder bistatic paths; such uncertainty requires
scenario-specific bundle coefficients and a multi-scenario master.
